\scanpage{75}
% archive sheet number: 74
% L1-6
\uline{Proposition} Soit $G$ un groupe algébrique \uline{connexe} sur un corps \uline{fini} $k$, $L$ une extension finie de $k$, de groupe de Galois $\Gamma$, $G^L$ le groupe obtenu à partir de $G$ par extension des scalaires à $L$ (il est encore connexe sur $L$) $\Gamma$ opère dans $G^L$ comme groupe d'automorphismes. \deleted{Soit} \deleted{\unsure{$\psi$ un}\ \illegible{1}} Alors $H^1(\Gamma,G^L)=\{e\}$. Plus précisément,
% slanted; alongside a vertical emphasis stroke, not a cancellation
\margin{\unsure{6|b}\quad\unsure{pas vrai!}}
% L7-11
soit $\operatorname{Aut}G^L$ le groupe des automorphismes du groupe algébrique $G^L$ \uline{sur $L$}, et soit $\psi$ un \deleted{\unsure{groupe}}\inserted{hom} croisé de $\Gamma$ dans $\operatorname{Aut}G^L$, soit $G^{L\psi}$ le groupe $G^L$ \deleted{\unsure{muni}} où $\Gamma$ opère par $\sigma(\gamma).g=\psi(\gamma).g^\gamma$ ; alors on a encore
\[
H^1(\Gamma,G^{L\psi})=\{e\}
\]

% L12-15
\uline{Démonstration} Soit $q=\text{nb d'éléments de }k$, et $n=[L:k]$.
$\Gamma$ est le groupe cyclique engendré par l'automorphisme $\unsure{f|u}\lambda=\lambda^q$, on désigne par $u$ l'automorphisme correspondant de $G^L$. Soit $\Omega$ une clôture algébrique de $L$, on désigne par $\widetilde u$ l'automorphisme $\deleted{\unsure{\widetilde f\lambda=\lambda^q\text{ de }\Omega}}$, \deleted{\unsure{et par $\widetilde u$ aussi}} \deleted{$\unsure{x\longrightarrow x^{\widetilde u}}$} \deleted{l'automorphisme} \deleted{$x\longrightarrow x^q$} de $G^\Omega$ \deleted{\unsure{défini ainsi}} défini \inserted{par l'automorphisme $\lambda\longrightarrow\lambda^q$ de $\Omega$} (noter que $u^n=1$, mais $\widetilde u^n\ne1$ ; l'ensemble des invariants de $\widetilde u^n$ dans $\Omega$ --- resp$^t$ $G^\Omega$ --- est $L$ --- resp. $G^L$ ---).
% L16-20
On désigne par $\widetilde\sigma$ l'automorphisme $x\longrightarrow\psi(f)x^{\widetilde u}$ de $G^\Omega$, induisant un automorphisme $\sigma$ de $G$, qui n'est autre que $\sigma(f)$. On a, puisque $\psi(\unsure{\gamma|f})$ est un \deleted{$\unsure{\psi\,\psi^{\widetilde u}\psi^{\widetilde u^2}\cdots\psi^{\widetilde u}}$} hom croisé de $\Gamma$ dans $\operatorname{Aut}G^L$
